% mazzi: 04
\chapter{La trasmissione sinaptica}

% ---------------------------------------------------------------------------
\section{La comunicazione tra cellule}

Le cellule comunicano tra loro in 4 diverse modalità.
\begin{enumerate}
  \item \textbf{nervosa}: segnale elettrico che percorre l'assone e determina, al terminale
    assonico, il rilascio di un \textit{neurotrasmettitore} sulla cellula bersaglio; il segnale
    elettrico diventa chimico e sulla cellula bersaglio torna elettrico
    $\rightarrow$ risposta;
  \item \textbf{endocrina}: segnale chimico; l'ormone è secreto nel sangue e raggiunge cellule
    bersaglio anche lontane;
  \item \textbf{neuroendocrina}: segnale elettrico in una cellula nervosa (p.\,es.\ ipotalamica)
    che rilascia un ormone nel sangue diretto alla cellula bersaglio;
  \item \textbf{paracrina e autocrina}: segnale chimico; l'ormone è rilasciato a breve distanza nel
    liquido extracellulare e agisce su cellule vicine (paracrina) o sulle stesse cellule secernenti
    (autocrina).
\end{enumerate}

\rubrica{Specificità del segnale ormonale}
Gli ormoni sono immessi nel sangue e raggiungono tutte le cellule dell'organismo, ma sono
riconosciuti solo dalle \textbf{cellule bersaglio}, dotate di recettori specifici per quell'ormone;
dove il recettore manca non c'è risposta. L'interazione ormone-recettore genera diverse risposte
cellulari:
\begin{itemize}
  \item sintesi di proteine;
  \item attivazione di proteine;
  \item apertura o chiusura di canali ionici.
\end{itemize}

% ---------------------------------------------------------------------------
\section{Le sinapsi}

\textbf{Sinapsi}: punto di contatto \textbf{funzionale e non fisico} tra cellule, dove avviene la
trasmissione del segnale da un elemento all'altro; le due cellule in comunicazione non si toccano
mai. È la modalità con cui comunicano le cellule eccitabili.

Si distinguono essenzialmente in \textbf{sinapsi elettriche} e \textbf{sinapsi chimiche}.

% ---------------------------------------------------------------------------
\section{Le sinapsi elettriche}

\textbf{Sinapsi elettrica}: unione delle membrane di due cellule che si accostano fino a una
distanza di 3--4\,nm, formando una \textbf{giunzione comunicante} o \textit{gap junction}. Non
utilizzano l'intermediazione di una sostanza chimica.
\begin{itemize}
  \item connessione a \textbf{bassa resistenza}: permette il passaggio diretto e passivo di ioni da
    una cellula all'altra, attraverso canali di membrana affacciati tra elemento pre- e
    postsinaptico;
  \item quando il neurone presinaptico si depolarizza, il $Na^+$ accumulato al suo interno passa
    direttamente alla cellula postsinaptica.
\end{itemize}

\rubrica{Sedi}
Presenti p.\,es.\ nei nuclei ipotalamici neuroendocrini, ma anche in cellule non nervose (miociti).

\rubrica{Caratteristiche della trasmissione}
\begin{itemize}
  \item avviene \textbf{senza alcun ritardo};
  \item avviene in \textbf{ambedue le direzioni} e dipende dal gradiente di concentrazione degli
    ioni;
  \item induce nell'elemento postsinaptico modificazioni dello \textbf{stesso segno} dell'elemento
    presinaptico: se la cellula presinaptica si depolarizza, si depolarizza anche la postsinaptica.
\end{itemize}

% ---------------------------------------------------------------------------
\section{Le sinapsi chimiche}

Sono di gran lunga le più diffuse. Vi si distinguono tre componenti.
\begin{enumerate}
  \item \textbf{elemento presinaptico}: contiene mitocondri, elementi tubulari e \textit{vescicole
    sinaptiche} contenenti il neurotrasmettitore. I mitocondri forniscono l'ATP necessario a
    impacchettare i neurotrasmettitori e a permetterne il rilascio;
  \item \textbf{spazio sinaptico}: 20--40\,nm di spessore (più ampio di quello delle sinapsi
    elettriche);
  \item \textbf{elemento postsinaptico}: la sua membrana presenta \textbf{recettori} specifici
    (proteine) cui si lega il neurotrasmettitore; è il legame neurotrasmettitore-recettore a
    produrre l'effetto biologico nella cellula postsinaptica.
\end{enumerate}

\rubrica{Cellule coinvolte e tipi di contatto}
Le sinapsi possono avvenire tra cellule nervose, tra cellule nervose e cellule muscolari o tra
cellule nervose e cellule ghiandolari. Secondo la sede del contatto sono di tipo:
\begin{itemize}
  \item \textbf{asso-dendritico}: assone della cellula presinaptica $\rightarrow$ dendriti della
    postsinaptica;
  \item \textbf{asso-somatico}: assone $\rightarrow$ corpo cellulare;
  \item \textbf{asso-assonico}: assone $\rightarrow$ assone di un'altra cellula.
\end{itemize}

\rubrica{Caratteristiche della trasmissione}
\begin{itemize}
  \item si verifica con un \textbf{ritardo} di circa 0,5\,ms, dovuto al rilascio del
    neurotrasmettitore e al suo attacco al recettore: tutti gli eventi a questo livello richiedono
    tempo;
  \item è \textbf{unidirezionale}: solo dall'elemento presinaptico a quello postsinaptico;
  \item può indurre nell'elemento postsinaptico eventi sia \textbf{eccitatori} --- EPSP
    (\textit{excitatory post-synaptic potential}) --- sia \textbf{inibitori} --- IPSP
    (\textit{inhibitory post-synaptic potential}) --- a seconda del canale ionico che si apre
    ($Na^+$, $K^+$, $Cl^-$).
\end{itemize}

\subsection{Sintesi e rilascio del neurotrasmettitore}

\begin{enumerate}
  \item \textbf{sintesi degli enzimi} nel corpo cellulare (RER, apparato del Golgi);
  \item \textbf{trasporto assonico lento degli enzimi} lungo i microtubuli dell'assone;
  \item \textbf{sintesi e immagazzinamento del neurotrasmettitore} nella terminazione presinaptica:
    a partire dai precursori e per opera degli enzimi, il neurotrasmettitore viene impacchettato
    nelle vescicole sinaptiche, che si ancorano al citoscheletro in attesa del momento del rilascio;
  \item \textbf{rilascio e diffusione del neurotrasmettitore} nello spazio sinaptico;
  \item \textbf{trasporto dei precursori} nella terminazione presinaptica.
\end{enumerate}

% ---------------------------------------------------------------------------
\section{La giunzione neuromuscolare}

\textbf{Giunzione neuromuscolare}: sinapsi tra un neurone e una cellula muscolare. Si attiva quando
il sistema nervoso centrale comanda una contrazione. Il neurotrasmettitore è unico:
l'\textbf{acetilcolina} (ACh).

\rubrica{Sequenza degli eventi}
\begin{enumerate}
  \item \textbf{potenziale d'azione} nella fibra nervosa: generato nel motoneurone delle corna
    grigie midollari, viaggia lungo l'assone dell'$\alpha$-motoneurone e depolarizza il terminale
    presinaptico;
  \item ingresso di \textbf{$Ca^{2+}$} nel terminale nervoso presinaptico: le vescicole sinaptiche
    si staccano dal citoscheletro, raggiungono la membrana e si fondono con essa;
  \item \textbf{rilascio dell'ACh} nello spazio sinaptico;
  \item \textbf{diffusione dell'ACh} verso la placca motrice;
  \item \textbf{legame ACh-recettore} sulla placca motrice: si aprono i canali per $Na^+$ e $K^+$;
    l'ingresso di $Na^+$ depolarizza la membrana;
  \item \textbf{depolarizzazione della placca motrice} $\rightarrow$ potenziale d'azione nel
    muscolo $\rightarrow$ inizio della contrazione;
  \item \textbf{degradazione dell'ACh}: esaurito il compito, il neurotrasmettitore si stacca dal
    recettore e torna libero nello spazio sinaptico, dove l'enzima \textbf{acetilcolinesterasi}
    (AChE) lo degrada in \textit{colina} e \textit{acetato}.
\end{enumerate}

\rubrica{Ciclo di sintesi e degradazione dell'acetilcolina}
\begin{itemize}
  \item \textbf{sintesi}: colina $+$ acetil-CoA $\rightarrow$ acetilcolina, per opera della
    \textit{colina acetiltransferasi}. L'acetil-CoA deriva dal metabolismo cellulare del glucosio
    (ciclo di Krebs);
  \item \textbf{degradazione}: acetilcolina $\rightarrow$ colina $+$ acetato, per opera
    dell'\textit{acetilcolinesterasi};
  \item \textbf{ricaptazione}: colina e acetato rientrano nel terminale nervoso, dove l'acetilcolina
    viene riformata e reimpacchettata nelle vescicole, pronta per una nuova trasmissione.
\end{itemize}

Negli altri neurotrasmettitori non esistono enzimi specifici di degradazione: il neurotrasmettitore
viene ripreso tal quale dal terminale presinaptico o dalle cellule della glia.

% ---------------------------------------------------------------------------
\section{I potenziali postsinaptici}

A seconda del recettore e del canale che si apre, la risposta postsinaptica è eccitatoria o
inibitoria.

\rubrica{Sinapsi eccitatorie}
Recettori ionotropi che aprono canali per ioni positivi concentrati all'esterno, che entrano nella
cellula e depolarizzano la membrana avvicinandola alla soglia $\rightarrow$ \textbf{EPSP}.
\begin{itemize}
  \item ingresso di \textbf{$Na^+$};
  \item ingresso di \textbf{$Ca^{2+}$}.
\end{itemize}

\rubrica{Sinapsi inibitorie}
Recettori ionotropi il cui effetto netto è un'\textbf{iperpolarizzazione}, che allontana la
membrana dalla soglia $\rightarrow$ \textbf{IPSP}.
\begin{itemize}
  \item ingresso di \textbf{$Cl^-$}: ione negativo, più concentrato all'esterno, entra nella
    cellula;
  \item uscita di \textbf{$K^+$}: ione positivo, più concentrato all'interno, tende a uscire.
\end{itemize}

\subsection{Propagazione dell'EPSP e formazione del potenziale d'azione}

\begin{enumerate}
  \item l'\textbf{EPSP} è generato dalla stimolazione del terminale presinaptico a livello dei
    dendriti;
  \item l'EPSP \textbf{decade poco} lungo il soma ed è in grado di raggiungere la soglia di
    attivazione e di generare un potenziale d'azione;
  \item il \textbf{potenziale d'azione} generato a livello del monticolo assonico è in grado di
    propagarsi lungo l'assone.
\end{enumerate}

La differenza dipende dalla densità dei canali del $Na^+$ voltaggio-dipendenti:
\begin{itemize}
  \item \textbf{dendriti e soma}: bassa densità $\rightarrow$ alta soglia di attivazione per il
    potenziale d'azione;
  \item \textbf{monticolo assonico}: alta densità $\rightarrow$ bassa soglia di attivazione.
\end{itemize}

Il potenziale d'azione è un fenomeno \textbf{tutto o nulla}: una volta formato si trasmette lungo
l'assone sempre uguale a se stesso, e al terminale è sicuramente in grado di far liberare il
neurotrasmettitore. I potenziali \textit{elettrotonici} (EPSP e IPSP), al contrario, si attenuano
propagandosi.

% ---------------------------------------------------------------------------
\section{La sommazione sinaptica}

Ogni cellula nervosa riceve migliaia di contatti sinaptici, sia eccitatori sia inibitori, dai
terminali assonici dei neuroni presinaptici sui propri dendriti e sul soma (tra i quali si
insinuano i processi delle cellule gliali). Tutti questi potenziali convergono sul
\textbf{monticolo assonico}, dove avviene la \textbf{sommazione}. Se ne distinguono due tipi.

\subsection{Sommazione spaziale}

Gli input sinaptici arrivano in \textbf{punti diversi} della cellula (e in momenti diversi) da
cellule presinaptiche differenti. Preso l'esempio di tre input A, B e C, ciascuno dei quali
depolarizza la membrana di 5\,mV, nessuno raggiunge da solo la soglia.
\begin{itemize}
  \item \textbf{solo EPSP}: se i tre potenziali arrivano al monticolo assonico prima che la cellula
    si sia ripolarizzata a $-70$\,mV, si sommano ($5 + 5 + 5 = 15$\,mV): da $-70$\,mV si raggiunge
    la soglia di $-55$\,mV $\rightarrow$ si aprono i canali del $Na^+$ voltaggio-dipendenti
    $\rightarrow$ \textbf{potenziale d'azione} (sommazione spaziale sopra soglia);
  \item \textbf{EPSP e IPSP}: se uno dei tre input è inibitorio (p.\,es.\ C, che apre i canali del
    $Cl^-$), la somma diventa $+5 +5 -5 = +10$\,mV: da $-70$\,mV si arriva a $-60$\,mV, la soglia
    non è raggiunta e si ha solo un EPSP, che si attenua lungo la propagazione (sommazione spaziale
    sotto soglia).
\end{itemize}

\subsection{Sommazione temporale}

Gli input arrivano \textbf{sempre dallo stesso terminale assonico}, cioè dalla stessa cellula
presinaptica, e sono quindi dello stesso segno: per il \textit{principio di Dale} un neurone è
capace di produrre un solo tipo di neurotrasmettitore, e troverà quindi recettori dello stesso
segno.
\begin{itemize}
  \item potenziali in \textbf{rapida successione}: si sommano $\rightarrow$ si raggiunge la soglia
    $\rightarrow$ potenziale d'azione;
  \item potenziali \textbf{distanziati nel tempo}: la membrana ha il tempo di tornare a $-70$\,mV,
    non si sommano $\rightarrow$ nessun potenziale d'azione.
\end{itemize}

% ---------------------------------------------------------------------------
\section{I neurotrasmettitori}

\textbf{Neurotrasmettitori}: sostanze chimiche di vario tipo, sintetizzate nel corpo cellulare del
neurone e riconosciute da recettori specifici sulla membrana postsinaptica. Si dividono in:
\begin{itemize}
  \item \textbf{molecole di basso peso molecolare}: acetilcolina, amine; e amminoacidi
    (glutammato, aspartato, acido $\gamma$-amminobutirrico o GABA, glicina);
  \item \textbf{peptidi}: oltre 100, formati per la maggior parte da catene di 3--30 amminoacidi.
\end{itemize}

\rubrica{Distribuzione}
\begin{itemize}
  \item \textbf{sistema nervoso somatico periferico}: il neurotrasmettitore per eccellenza è
    l'\textbf{acetilcolina}, rilasciata a livello della giunzione neuromuscolare;
  \item \textbf{sistema nervoso centrale}: identificati molti neurotrasmettitori, tra cui la
    \textbf{dopamina} (famiglia delle \textit{catecolammine}, con noradrenalina e adrenalina), la
    \textbf{serotonina}, il \textbf{glutammato}, la \textbf{glicina}, il \textbf{GABA}.
\end{itemize}

\rubrica{Neurotrasmettitori occasionali}
Alcune sostanze si comportano da neurotrasmettitori solo in determinate situazioni: diverse purine
(ATP, adenosina, AMP), l'ossido nitrico (NO), il monossido di carbonio (CO).

\rubrica{Corredo recettoriale della cellula postsinaptica}
Normalmente una cellula possiede recettori differenti, non un solo tipo; l'unica eccezione è la
cellula muscolare, che ha solo recettori per l'acetilcolina. Inoltre uno stesso neurotrasmettitore
(p.\,es.\ il glutammato) può disporre sulla cellula postsinaptica di più tipi di recettore, con
effetti differenti.

\subsection{Sintesi e degradazione delle catecolammine}

\rubrica{Via di sintesi}
$$\text{Tirosina} \rightarrow \text{L-Dopa} \rightarrow \text{Dopamina} \rightarrow
\text{Noradrenalina} \rightarrow \text{Adrenalina}$$
\begin{itemize}
  \item tirosina $\rightarrow$ L-Dopa: \textit{tirosina-idrossilasi};
  \item L-Dopa $\rightarrow$ dopamina: \textit{dopa-decarbossilasi} (neurone dopaminergico);
  \item dopamina $\rightarrow$ noradrenalina: \textit{dopamina-$\beta$-idrossilasi} (neuroni
    adrenergici);
  \item noradrenalina $\rightarrow$ adrenalina:
    \textit{feniletanolamina-N-metiltransferasi} (midollare del surrene).
\end{itemize}

\rubrica{Vie di degradazione}
Operano due enzimi, le \textbf{monoamminossidasi} (MAO) e la \textbf{catecol-O-metiltransferasi}
(COMT), da soli o in combinazione.

\begin{table}[htbp]
  \centering
  \small
  \renewcommand{\arraystretch}{1.3}
  \begin{tabularx}{\textwidth}{|l|X|X|X|}
    \hline
    \textbf{Catecolammina} & \textbf{MAO} & \textbf{COMT} & \textbf{MAO $+$ COMT} \\
    \hline
    Dopamina & acido diidrossifenilacetico & 3-metossitiramina & acido omovanillico (HVA) \\
    \hline
    Noradrenalina & acido diidrossimandelico & normetanefrina
      & acido 3-metossi-4-idrossimandelico (VMA) \\
    \hline
    Adrenalina & acido diidrossimandelico & metanefrina
      & acido 3-metossi-4-idrossimandelico (VMA) \\
    \hline
  \end{tabularx}
  \caption{Metaboliti di degradazione delle catecolammine secondo l'enzima coinvolto.}
  \label{tab:degradazione-catecolammine}
\end{table}

% ---------------------------------------------------------------------------
\section{Effetti del legame neurotrasmettitore-recettore}

I recettori postsinaptici si dividono in due grandi categorie.

\subsection{Recettori ionotropi}

\textbf{Recettori ionotropi}: canali ionici regolati chimicamente; il legame del neurotrasmettitore
agisce \textbf{direttamente} sul canale, con azione immediata $\rightarrow$ potenziale sinaptico
\textbf{rapido e di breve durata}.
\begin{itemize}
  \item \textbf{canali che si aprono}:
    \begin{itemize}
      \item aumentato ingresso di $Na^+$ $\rightarrow$ \textbf{EPSP}: depolarizzazione, eccitazione;
      \item aumentata uscita di $K^+$ o ingresso di $Cl^-$ $\rightarrow$ \textbf{IPSP}:
        iperpolarizzazione, inibizione;
    \end{itemize}
  \item \textbf{canali che si chiudono}:
    \begin{itemize}
      \item diminuito ingresso di $Na^+$ $\rightarrow$ \textbf{IPSP}: iperpolarizzazione,
        inibizione;
      \item diminuita uscita di $K^+$ $\rightarrow$ \textbf{EPSP}: depolarizzazione, eccitazione.
    \end{itemize}
\end{itemize}

\subsection{Recettori metabotropi}

\textbf{Recettori metabotropi}: recettori legati a una \textbf{proteina G}; l'azione sul canale non
è immediata ma ritardata, perché mediata dall'attivazione della via del \textbf{secondo
messaggero} $\rightarrow$ potenziali sinaptici \textbf{lenti} ed effetti a lungo termine. Il nome
risale all'ipotesi, poi superata, che l'evento intermedio fosse un processo metabolico.

Effetti possibili:
\begin{itemize}
  \item \textbf{modifica dello stato di apertura di canali ionici}: apertura o chiusura (a volte è
    necessario interrompere un'azione);
  \item \textbf{modifica di proteine esistenti o regolazione della sintesi di nuove proteine}, fino
    a interferire con la trascrizione genica $\rightarrow$ \textbf{risposta intracellulare
    coordinata}.
\end{itemize}

% ---------------------------------------------------------------------------
\section{Il ruolo modulatore delle sinapsi chimiche}

Le sinapsi chimiche modulano tutte le nostre attività modificando il segnale in ingresso:
\begin{enumerate}
  \item \textbf{amplificazione} del segnale;
  \item \textbf{inversione} del segnale: tramite le sinapsi inibitorie si limita un'eccitazione
    eccessiva;
  \item \textbf{prolungamento temporale} dell'azione;
  \item \textbf{fenomeni di plasticità}: il cervello si rimodella nel corso della vita, e il primo
    livello di rimodellamento è proprio la trasmissione sinaptica. L'apprendimento, anche di un
    gesto motorio, comporta la formazione di nuove spine dendritiche e di nuovi recettori di
    membrana $\rightarrow$ risposta più pronta allo stesso stimolo. La prima fase è una
    modificazione \textit{funzionale}, cui segue una modificazione \textit{strutturale}.
\end{enumerate}
